% Labb 3, parts 18 to 23, from labs/lab3/d_io_and_subroutines.md. Each part is a \labtask, and its
% Mål, Bakgrund, Programmet, Uppgifter and Frågor are \task headings. The programs given in parts
% 20 (with its two constants left to fill in) and 23 (with its deliberate fault) are the
% handout's. One deliberate difference: the last comment of the program in part 23 says to copy
% delay_ms from part 22, where the guide says L10 Appendix C.4, a name the book does not use.

\section{Del 18--23: Text, 16 bitar, inporten och subrutiner}
\label{lab3:sec:d}
Delarna hör till kapitel~\ref{c10:ch}. Läs \secref{c10:app:a} före del 18--20, \secref{c10:app:b}
före del~21 och \secref{c10:app:c} före del 22--23. Instruktionerna finns på referensbladet i
bilaga~\ref{app:avr}.

Del 21 krävs för godkänt. Del 18--20 och 22--23 krävs för VG, och del~22 behövs dessutom för
slutuppgiften i kapitel~\ref{c11:ch}, så gör den även om ni inte siktar på VG.

\labtask{Del 18}{ASCII-koder för hexadecimala siffertecken}\phantomsection\label{lab3:d18}
\task{Mål}
Göra om värden till de tecken som skriver ut dem, och tillbaka, och lägga text i minnet.

\task{Bakgrund}
ASCII-tabellen och luckan på sju mellan \code{'9'} och \code{'A'} finns i
\secref{c10:a1}--\ref{c10:a2}. En byte blir två tecken med \code{swap} och \code{andi} enligt
\secref{c10:a3}, och \code{sts} beskrivs i \secref{c10:a5}. Kom ihåg att \code{subi r17, -0x30}
adderar \code{0x30}.

\task{Uppgifter}
\begin{enumerate}
  \item \textbf{En siffra blir ett tecken.} Ladda \code{r17} med 7 och gör om det till tecknet
    \code{'7'} med en enda instruktion. \emph{Förutsäg} \code{r17} hexadecimalt innan ni stegar.
  \item \textbf{En hexadecimal siffra blir ett tecken.} Ladda \code{r18} med ett värde 0--15,
    kopiera det till \code{r19} och gör om \code{r19} till tecknet för siffran, med \code{cpi},
    \code{brlo} och två \code{subi}. Fyll i tabellen \textbf{innan} ni kör programmet, och
    kontrollera sedan varje rad genom att byta värdet i \code{ldi}:

    \begin{filltable}{@{}p{5em}YY@{}}
      \thead{\code{r18}} & \thead{\code{r19}, förutsagt} & \thead{\code{r19}, i simulatorn} \\
      \midrule
      \code{0x00} & & \fillrule
      \code{0x09} & & \fillrule
      \code{0x0A} & & \fillrule
      \code{0x0B} & & \fillrule
      \code{0x0F} & & \\
    \end{filltable}
  \item \textbf{En byte blir text.} Ladda \code{r16} med \code{0xA7}. Gör om den höga nibblen till
    ett tecken i \code{r20} och den låga till ett tecken i \code{r21}, och lägg dem på adress
    \code{0x0100} och \code{0x0101} med \code{sts}. Öppna fönstret \textbf{Memory}, välj
    \textbf{data IRAM} och gå till adress \code{0x0100}. Står texten \code{A7} i ASCII-kolumnen
    till höger? Prova också \code{0x3C}, \code{0xFF} och \code{0x00}.
  \item \textbf{Tillbaka från tecken till värde.} Ladda \code{r22} med \code{'F'} och gör om det
    till värdet 15, enligt \secref{c10:a4}. Prova också \code{'0'}, \code{'9'} och \code{'A'}.
\end{enumerate}

\task{Frågor}
\begin{itemize}
  \item Varför måste den höga nibblen flyttas ned med \code{swap} innan den kan göras om till ett
    tecken?
  \item Vad blir tecknet om man glömmer \code{andi r20, 0x0F} efter \code{swap}, för byten
    \code{0xA7}?
\end{itemize}

\redovisa{uppgift 3 och 4 i simulatorn, med texten synlig i Memory-fönstret.}

\labtask{Del 19}{16-bitarsoperationer}\phantomsection\label{lab3:d19}
\task{Mål}
Räkna med tal som inte får plats i en byte: addera, subtrahera, räkna och jämföra 16-bitarstal.

\task{Bakgrund}
Registerpar och \code{low()}/\code{high()} finns i \secref{c10:a6}, addition och subtraktion med
minnessiffra i \secref{c10:a7}, \code{adiw}, \code{sbiw} och \code{movw} i \secref{c10:a8} och
jämförelsen med \code{cpi}/\code{cpc} i \secref{c10:a9}.

\task{Uppgifter}
\begin{enumerate}
  \item \textbf{Addition.} Beräkna 1000 + 2000 med \code{r25:r24} och \code{r23:r22}, med
    \code{add} och \code{adc}. Kopiera resultatet till \code{r19:r18} med \code{movw}.

    \emph{Förutsäg} först, med additionen uppställd för hand: vad blir \code{r24} och flaggan
    \code{C} efter \code{add}, och vad blir \code{r25} efter \code{adc}? Stega sedan och jämför i
    Processor Status.
  \item \textbf{Subtraktion.} Beräkna 3000 -- 1234 med \code{sub} och \code{sbc}, och lägg
    resultatet i \code{r21:r20}. \emph{Förutsäg} båda bytena hexadecimalt. Blir det ett lån från
    den låga byten?
  \item \textbf{En 16-bitars räknare.} Nollställ \code{r25:r24} och räkna upp det med \code{adiw}
    tills det är 1000, med \code{cpi} och \code{cpc} i loopen. Sätt en etikett direkt efter loopen.

    \emph{Förutsäg} hur många cykler loopen tar, med samma metod som i kapitel~\ref{c9:ch}. Mät
    sedan med brytpunkter före och efter loopen och fönstret Processor Status
    (\secref{studio:4-4}). Stämmer siffrorna? Om inte: vilket varv räknade ni fel på?
  \item \textbf{När 16 bitar inte räcker.} Ladda \code{r27:r26} med \code{0xFFFF} och addera 1 med
    \code{adiw}. \emph{Förutsäg} resultatet och flaggorna \code{C} och \code{Z}.
\end{enumerate}

\task{Frågor}
\begin{itemize}
  \item Varför måste de låga bytena adderas först?
  \item Varför behövs registret med \code{high(1000)} i uppgift~3, när den låga byten kan jämföras
    direkt med \code{cpi}?
\end{itemize}

\redovisa{uppgift 1--4, och er förutsägelse och mätning i uppgift~3.}

\labtask{Del 20}{Lång tidsfördröjning}\phantomsection\label{lab3:d20}
\task{Mål}
Konstruera en fördröjning på en halv sekund, exakt på cykeln, och få en lysdiod att blinka en gång
per sekund.

\task{Bakgrund}
Metoden, en 16-bitars inre loop i en yttre 8-bitars loop, och formeln \code{M(4N + 4)} finns i
\secref{c10:a10}. Exemplet där blinkar med en sekunds halvperiod; här ska det gå dubbelt så fort.

\task{Programmet}
Fyll i de två konstanterna.

\begin{asmcode}
.include "m328Pdef.inc"

.equ OUTER_TURNS = ?                ; Work these two out before you run anything.
.equ INNER_TURNS = ?

.org 0x0000
    rjmp main

main:
    sbi DDRB, 5                     ; PB5, Arduino pin 13, is an output.
    ldi r17, 0b00100000             ; The mask that selects bit 5.

loop:
    in r16, PORTB                   ; Toggle PB5.
    eor r16, r17
    out PORTB, r16

delay_begin:
    ldi r20, OUTER_TURNS
delay_outer:
    ldi r24, low(INNER_TURNS)
    ldi r25, high(INNER_TURNS)
delay_inner:
    sbiw r24, 1
    brne delay_inner
    dec r20
    brne delay_outer
delay_end:
    rjmp loop
\end{asmcode}

\task{Uppgifter}
\begin{enumerate}
  \item \textbf{Räkna först.} Hur många cykler är en halv sekund vid 16~MHz? Välj
    \code{INNER_TURNS} så att ett varv i den yttre loopen tar 200\,000 cykler, och räkna ut
    \code{OUTER_TURNS}. Skriv upp uträkningen.
  \item \textbf{Mät.} Sätt stoppurets frekvens till 16~MHz, sätt en brytpunkt på
    \code{delay_begin} och en på \code{delay_end}, och mät tiden mellan dem. \emph{Förutsäg}
    stoppurets värde innan ni kör.
  \item \textbf{Hela perioden.} Hur lång är en hel blinkperiod, från en tändning till nästa, räknat
    i cykler? Varför är den inte exakt en sekund? Hur mycket fel är den, i mikrosekunder?
  \item \textbf{En annan tid.} Ändra till en fjärdedels sekund. Vilken av de två konstanterna ändrar
    ni, och till vad?
\end{enumerate}

\begin{aside}
\textbf{Tips.} Simulatorn behöver en stund för att köra åtta miljoner cykler. Använd
\textbf{Continue} (\textbf{F5}) mellan brytpunkterna, inte \textbf{Step Into}.
\end{aside}

\task{Frågor}
\begin{itemize}
  \item Varför är det enklare att välja ett \enquote{jämnt} antal cykler för den yttre loopen först,
    och räkna ut antalet varv sedan?
  \item Vad händer om \code{OUTER_TURNS} sätts till 0?
\end{itemize}

\redovisa{uträkningen och mätningen.}

\labtask{Del 21}{Inporten PORTD (PIND)}\phantomsection\label{lab3:d21}
\task{Mål}
Läsa knappar på port D och styra lysdioder på port B utifrån dem, och bygga logiska funktioner i
mjukvara.

\task{Bakgrund}
Pull-up-motståndet och varför en nedtryckt knapp läses som 0 finns i \secref{c10:b2}, att läsa
hela porten i \secref{c10:b3} och \code{sbic}/\code{sbis} i \secref{c10:b4}. Så simuleras en
knapp: \secref{studio:4-5}.

I hela delen sitter knapparna mellan stiftet och jord, och de inbyggda pull-uparna används. En
\textbf{tom} ruta i \code{PIND} är alltså en \textbf{nedtryckt} knapp.

\task{Uppgifter}
\begin{enumerate}
  \item \textbf{Knapparna på lysdioderna.} Skriv av programmet \code{buttons_to_leds.asm} från
    \secref{c10:b3}. Kör till loopen, stoppa, och ändra \code{PIND} för hand. \emph{Förutsäg}
    \code{PORTB} för varje rad innan ni stegar:

    \begin{filltable}{@{}p{6.4em}YYY@{}}
      \thead{\code{PIND}} & \thead{Nedtryckta knappar} & \thead{\code{PORTB}, förutsagt} &
        \thead{\code{PORTB}, i simulatorn} \\
      \midrule
      \code{0b11111111} & & & \fillrule
      \code{0b11111011} & & & \fillrule
      \code{0b11110011} & & & \fillrule
      \code{0b01111110} & & & \\
    \end{filltable}
  \item \textbf{En knapp, en lysdiod.} Skriv ett program där lysdioden på PB5 lyser så länge
    knappen på PD2 hålls nedtryckt, med \code{sbic} eller \code{sbis} och utan att läsa hela porten.
  \item \textbf{AND och OR i mjukvara.} Två knappar, på PD2 och PD3. Lysdioden på PB0 ska lysa när
    \textbf{båda} är nedtryckta, och lysdioden på PB1 när \textbf{minst en} är det. Det är samma två
    funktioner som seriekopplingen och parallellkopplingen i Labb~1 (bilaga~\ref{app:lab1}).

    Rita en flödesplan först. Fyll sedan i sanningstabellen, med förutsägelsen före simuleringen:

    \begin{filltable}{@{}p{5.6em}p{5.6em}YY@{}}
      \thead{PD2} & \thead{PD3} & \thead{PB0 (AND)} & \thead{PB1 (OR)} \\
      \midrule
      släppt & släppt & & \fillrule
      släppt & nedtryckt & & \fillrule
      nedtryckt & släppt & & \fillrule
      nedtryckt & nedtryckt & & \\
    \end{filltable}
  \item \textbf{Extra.} Lägg till lysdioden på PB2, som ska lysa när \textbf{exakt en} av knapparna
    är nedtryckt. Vilken grind från kapitel~\ref{c2:ch} är det? Vilken koppling i Labb~1 gjorde
    samma sak?
\end{enumerate}

\task{Frågor}
\begin{itemize}
  \item Vad läser programmet på en ingång utan pull-up, när knappen är släppt?
  \item I uppgift~3 behövs \code{com} och \code{andi} innan jämförelsen. Vad gör var och en av dem,
    och vad går fel om man hoppar över \code{andi}?
  \item Varför är det enklare att ändra funktionen i uppgift~3 än i Labb~1?
\end{itemize}

\redovisa{uppgift~3, med sanningstabellen, i simulatorn.}

\labtask{Del 22}{Subrutiner}\phantomsection\label{lab3:d22}
\task{Mål}
Flytta en fördröjning till en subrutin, se returadressen på stacken, och skriva en subrutin som tar
en parameter.

\task{Bakgrund}
\code{rcall}, \code{ret}, stacken och stackpekaren finns i \secref{c10:c2}--\ref{c10:c3},
subrutinen \code{delay_ms} med parameter i \secref{c10:c4}, och en subrutin som anropar en annan i
\secref{c10:c7}.

Från och med den här delen ska varje program som använder \code{rcall} börja med att initiera
stackpekaren, med de fyra raderna i \secref{c10:c3}.

\task{Uppgifter}
\begin{enumerate}
  \item \textbf{Fördröjningen blir en subrutin.} Ta programmet från del~20. Flytta fördröjningen,
    från \code{delay_begin} till och med \code{brne delay_outer}, till en subrutin
    \code{delay_500ms} som slutar med \code{ret}, och anropa den med \code{rcall} där fördröjningen
    stod. Lägg till initieringen av \code{SP}. Programmet ska blinka precis som förut.
  \item \textbf{Returadressen.} \emph{Förutsäg} \code{SP} före \code{rcall}, i början av
    subrutinen och efter \code{ret}. Leta upp adressen till instruktionen efter \code{rcall} i
    listfilen (\code{.lss}, \secref{c7:b7}). Stega in i subrutinen med \textbf{F11} och läs de två
    översta bytena i stacken i fönstret Memory (\textbf{data IRAM}, adress
    \code{0x08FE}--\code{0x08FF}). Är det returadressen? Vilken byte ligger var?
  \item \textbf{En parameter.} Skriv av \code{delay_ms} från \secref{c10:c4}. Skriv om
    \code{delay_500ms} så att den laddar \code{r24} med 250 och anropar \code{delay_ms} två gånger.
    Glöm inte att \code{delay_500ms} själv ändrar \code{r24}, och alltså ska spara det.
  \item \textbf{Järnvägsövergången.} Ändra huvudprogrammet så att lysdioderna på PB0 och PB1
    blinkar växelvis, en halv sekund var, som vid en järnvägsövergång.
  \item \textbf{Mät ett anrop.} \emph{Förutsäg} hur många cykler ett anrop av \code{delay_500ms}
    tar, från \code{rcall} till instruktionen efter, med formeln för \code{delay_ms} och cyklerna
    för \code{rcall}, \code{push}, \code{ldi}, \code{pop} och \code{ret}. Mät sedan med
    stoppuret.
  \item \textbf{Den djupaste punkten.} Stoppa programmet mitt i \code{delay_ms}, när den har
    anropats från \code{delay_500ms}. \emph{Förutsäg} \code{SP} först, och kontrollera. Hur många
    byte ligger på stacken, och vad är var och en?
\end{enumerate}

\task{Frågor}
\begin{itemize}
  \item Vad skulle hända om \code{delay_ms} inte sparade \code{r24}?
  \item \code{rcall} kostar 3 cykler och \code{ret} 4. Är det ett argument mot att använda
    subrutiner?
\end{itemize}

\redovisa{uppgift 2, 4 och 5.}

\labtask{Del 23}{Använda stacken som lagringsplats}\phantomsection\label{lab3:d23}
\task{Mål}
Använda \code{push} och \code{pop} för att spara register, förstå varför ordningen spelar roll, och
se vad som händer när stacken inte stämmer.

\task{Bakgrund}
\code{push}, \code{pop} och sist in, först ut finns i \secref{c10:c5}, och kursens regel om att
spara det man ändrar i \secref{c10:c6}.

\task{Programmet}
Huvudprogrammet räknar i \code{r16} hur många gånger det har anropat \code{flash_all}, och slutar
efter tre gånger. Subrutinen \code{flash_all} tänder alla lysdioder på port B i 100~ms och släcker
dem i 100~ms, och använder själv \code{r16} till något helt annat. Den saknar \code{push} och
\code{pop}, med flit.

\begin{asmcode}
.include "m328Pdef.inc"

.org 0x0000
    rjmp main

main:
    ldi r16, high(RAMEND)
    out SPH, r16
    ldi r16, low(RAMEND)
    out SPL, r16

    ldi r16, 0xFF
    out DDRB, r16                   ; Port B: all outputs.

    clr r16                         ; r16 counts the flashes.
loop:
    inc r16
    rcall flash_all
    cpi r16, 3
    brne loop
    out PORTB, r16                  ; Show the count.
end:
    rjmp end

flash_all:
    ldi r16, 0xFF                   ; All LEDs on ...
    out PORTB, r16
    ldi r24, 100
    rcall delay_ms                  ; ... for 100 ms ...
    clr r16                         ; ... and off ...
    out PORTB, r16
    ldi r24, 100
    rcall delay_ms                  ; ... for 100 ms.
    ret

; Copy delay_ms from part 22 here.
\end{asmcode}

\task{Uppgifter}
\begin{enumerate}
  \item \textbf{Byt plats med stacken.} Ladda \code{r20} med \code{0x55} och \code{r21} med
    \code{0xAA} i ett eget litet program. Lägg dem på stacken med \code{push r20}, \code{push r21},
    och hämta dem med \code{pop r20}, \code{pop r21}. \emph{Förutsäg} värdena i \code{r20} och
    \code{r21} efteråt, och \code{SP} efter varje instruktion.
  \item \textbf{Felet.} Kör programmet ovan. \emph{Förutsäg} först vad det gör. Hur många gånger
    blinkar det? Stega igenom ett varv i loopen och följ \code{r16}. Förklara vad som går fel.
  \item \textbf{Rättningen.} Lägg till \code{push} och \code{pop} i \code{flash_all} så att
    subrutinen följer kursens regel. Vilka register måste sparas? Programmet ska nu blinka tre
    gånger och sedan visa talet 3 på lysdioderna.
  \item \textbf{Djupet.} Stoppa programmet mitt i den första \code{delay_ms}. \emph{Förutsäg}
    \code{SP}, och kontrollera. Hur många byte ligger på stacken, och vad är var och en?
  \item \textbf{Fel ordning.} Byt plats på de två \code{pop} i \code{flash_all}, så att de inte
    längre står i omvänd ordning mot \code{push}. Vad händer? Förklara med stacken.
  \item \textbf{En pop för lite.} Ta bort en av \code{pop}-instruktionerna. Stega fram till
    \code{ret} i \code{flash_all} och läs \code{SP}. Stega en gång till: vart hoppar programmet?
    Förklara med stacken.
\end{enumerate}

\task{Frågor}
\begin{itemize}
  \item Varför måste \code{pop} göras i omvänd ordning mot \code{push}?
  \item Felet i uppgift~6 syns tydligare än felet i uppgift~5. Varför är det i själva verket en
    fördel?
\end{itemize}

\redovisa{uppgift 2, 3 och 4.}
